\subparagraph{Open problems.}

We now collect in physical terms the questions left open along the way, saying for each what is asked, why it matters for the construction and what would settle it.

The first is the kernel of the quartic form at higher level.
The auxiliary field lives on the quotient of the $2$-string states by the kernel of $\vertex{4}$, and the construction on the quotient assumes, from the quintic order on, either that the kernel is empty or that the quintic and sextic vertices vanish on it.
At level zero, for the tachyon pairs at every real pair momentum, the kernel is empty with Euclidean relative momenta, as in \cref{p:sec:worldsheet}, for the vertex family of~\cite{Erbin:2026:SL2CVerticesClosed} at the sections considered there with $\beta > \sqrt{2}$.
The pair kernel is then strictly definite on the even functions of the relative momentum.
At zero momentum the level-$(1, 1)$ block on the tachyon and the ghost dilaton is invertible (\cref{p:sec:solutions}).
Above level zero, the kernel is absent at generic momentum under the Fredholm property and the physical exclusion stated above.
  The question is whether it is absent at every real momentum.
  None of the mechanisms known to force zeros of string vertices (the dilaton theorem with the ghost-number counting, BRST-exact strands, the universality of the couplings at zero momentum), nor the worldsheet twist, produces a null vector of the quartic form at the levels and momenta where its entries are known.
  However, these mechanisms are not exhaustive.
  Nor can a positivity argument exclude a real zero: for the same vertex family, the level-$(1, 1)$ block is indefinite near zero momentum, so isolated real zeros are not excluded.
  The first block to compute is the level-$(1, 1)$ block at nonzero momentum, of which~\cite{Erbin:2026:SL2CVerticesClosed} now provides the first entries.

The question matters beyond consistency.
  At a real zero of the quartic form with a single null vector, the kernel condition of the quintic order reduces to the vanishing of a single pairing of the null vector with the quintic source at that momentum, a property of $\vertex{5}$ that nothing requires.
  If this pairing does not vanish, the composite vertex of \cref{p:sec:hs:quintic} has a pole in that momentum.
The reduced sextic vertex, $\vertex{6}$ minus six fifths of the composite vertex, then carries the opposite pole.
  The pole cancels in the combination $\dvertex{6} + \tfrac{6}{5} \, \vertex{5 \cdot 5} = \vertex{6}$.
  The pole is a propagating state in the $\Sigma_2$ channel, and it makes the reduced vertices singular there.
A real zero above level zero would be accidental (\cref{p:sec:hs:quartic}).
  We expect the pairing not to vanish on it, so that the multiplier on the kernel directions takes over there.
A computation of the momentum-dependent blocks of $\vertex{4}$ and $\vertex{5}$ at the first levels, from an explicit vertex family, would settle whether the quotient suffices or a multiplier on the kernel has to be used.

The second is the six-point exchange and the covering of moduli space.
The composite vertex $\vertex{5 \cdot 5}$, two quintic vertices joined by a $\Sigma_2$ line, has the dimension of $\mathcal M_{0,6}$.
Whether $\tfrac{6}{5} \vertex{5 \cdot 5}$ is the integral over a sub-region of $\mathcal V_{0,6}$, covered once, is open (\cref{p:sec:worldsheet}).
We conjecture that it is, with $\dvertex{6}$ the integral over the complement, so that the two cover the vertex region once.
  One route to a proof reads the composite vertex as glued surfaces, which requires the inverse of the quartic form to act as a gluing (\cref{p:sec:worldsheet}).
  For the explicit section of~\cite{Erbin:2026:SL2CVerticesClosed}, the inverse, summed over levels, does not act as a gluing with a Möbius cut coordinate, within the classes of gluings and of pair states specified there.
  Gluings with non-Möbius cut coordinates remain open.
At fixed level and genus zero, the results of \cref{p:sec:worldsheet} constrain the answer through the poles of the composite vertex in the pair momentum.
  A pole away from the plumbing lattice (the poles of the plumbing fixture) would rule out a compact sub-region with momentum-independent boundary and the standard measure, and the momentum-dependent regions of sector type considered there.
  We have neither established nor excluded such a pole.
  With a pole off the mass shells, on the plumbing lattice or not, or with no pole at all, the composite vertex is not a union of Feynman regions either, modulo the standard factorization theorem.
  The case of poles on the mass shells only is not settled.
  With a pole off the plumbing lattice, the object that remains is a momentum-weighted density on a region, such as a cusp-shaped region pinched at a boundary point, which can carry it.
  Whether such a region reproduces the composite vertex is open.

The question matters because it decides whether the reduced sextic vertex is itself a region of moduli space or an algebraic combination only, that is, what the exchange of $\Sigma_2$ corresponds to on moduli space at six points.
At fixed level, whether the composite vertex has a pole off the plumbing lattice is decided by computing the complex zero locus of the determinant of the quartic form in the pair momentum, with the residue at each zero.
If there are poles, matching their residues against a cusp-shaped region would decide at that level whether such a region reproduces the composite vertex; if there are none, the sub-region question stays open.
This computation needs the momentum-dependent blocks of the previous paragraph.
Even then, the fixed-level composite vertex is not the truncation of the full one, which inverts the quartic form on all $2$-string states at fixed external level.
A zero found at a fixed level is evidence only if it persists as the level grows.

The third is the level-zero zero.
At complex pair momentum, where a pole of the composite vertex off the plumbing lattice would sit, if there is one, the determinant of the discretized level-zero kernel vanishes at a point off the plumbing lattice.
This zero is stable under the discretization (\cref{p:sec:worldsheet}).
Whether the kernel itself has a square-integrable null function there is open.
The null vector of the discretized kernel is not resolved at the precision reached, and deciding it needs the kernel's entries to a precision well beyond it.
If a null function exists, its coupling to the quintic vertex decides at level zero whether the composite vertex has the pole of the previous paragraph.
As long as this is open, the pole hypothesis is neither realized nor excluded at level zero.
Whether a level-zero zero survives in the full composite vertex is the question of the level limit above.

The fourth is the conclusion of the mean-field approximation.
In the quantum theory the exact equation of the auxiliary field at quartic order says that $\Sigma_2$ is, formally, the pair correlator of the string field in its own background (\cref{p:sec:solutions}).
The Gaussian approximation to it gives a gap equation, as in the Hubbard--Stratonovich derivation of the BCS gap equation.
In the level-$(1, 1)$ instance at zero momentum, with the cubic vertex dropped, we solved that gap equation completely.
At the values of the coupling we examined, no root is a minimum of the potential on the real slice.
  Among the real roots, the only local convergent saddles on an admissible contour of the rotated $\Sigma_2$ integral are the one-loop displacements of the classical point.
A mean-field treatment thus gives back the classical structure and no new vacuum.

This computation has several limits.
  The mean field keeps the direct channel, one third of the one-loop self-energy (Hartree and Fock terms together), and drops the diagrams with an internal $\Sigma_2$ line, which no small parameter suppresses.
The obstacles to a BCS-type condensate are arguments, not theorems.
The instance has no loop momentum, no ghost sector in the trace and no cubic vertex.
We expect its conclusion to persist with loop momentum, but we have not established it.
  We have not examined the ghost sector and the cubic vertex.
The finite-dimensional models of \cref{p:app:finite-models} have no momenta either, and the mass shells and the on-shell behavior are outside them.
The gap equation with loop momentum, the ghost sector and the cubic vertex restored would settle whether the conclusion of the mean field persists.
The diagrams the mean field drops remain beyond this gap equation.

The fifth is the quantum master equation, which we do not address in this paper (\cref{p:sec:gauge}).
The construction is classical: the master action of the cubic theory is that of the original theory plus a trivial pair, through order $\gs^4$.
We do not include the loops of the auxiliary field, starting with the $\Sigma_2$ tadpole on the quintic vertex, nor the loop vertices of the original theory, which do not enter the classical construction.
  In this tadpole, the two $\Sigma_2$ legs of the coupling $\vertex{5}(\Sigma_2^2, \Psi)$ are joined by a $\Sigma_2$ propagator.
  Nor do we treat the Batalin--Vilkovisky structure of the multiplier form, whose multipliers need their own antifields.
  Beyond sextic order, we have not constructed the flat representative for the pair of $\Sigma_3$.
The question matters because the auxiliary field is a $2$-string field with its own propagator.
Whether its loops are absorbed in the loop vertices of the original theory, or generate vertices of their own, is what the quantum master equation would decide.
Settling it needs the master action with the loop vertices of the original theory and the auxiliary loops, and its master equation at one loop, with the trivial pair kept.

The sixth is the single-strand route, a different cubic theory suggested by the circle picture.
It has one asymmetric cubic vertex with a special puncture and an auxiliary field with a single string leg.
The quartic vertex would be produced by the exchange of that field after a change of its kinetic term, as in the open-string stub~\cite{Erbin:2024:OpenStringStub}.
We have established two things, the counting and the covering (\cref{p:sec:worldsheet}).
At four points, with the reversed plumbing range, every ghost and moduli count passes with sphere topology once the local discs of the asymmetric vertex extend beyond the unit circle.
The plumbing family covers the interior of $\mathcal V_{0,4}$ exactly once when the stub length varies along a seam symmetric in the three channels, the seam family of \cref{p:sec:worldsheet}, and never for a constant stub length.
This is numerical, at the parameters sampled, and the seam is under-determined.

Off-shell gauge invariance is not established: the local coordinates transported across the seam agree to first order only, and the strictly cubic single-strand theory is expected, though not shown, to fail gauge invariance off-shell at second order in the seam data on non-primary external states.
The question matters because such a theory would be cubic with a single string field and a single auxiliary field.
Its auxiliary channel would be a gluing region, which the inverse of the quartic form is not, in the class of kernels of \cref{p:sec:worldsheet}.
A check of the off-shell identity at second order, and of the embedding condition under which a quartic term on the collar of the seam restores it, would settle it.

The last is the $2$-string gauge parameter.
We have asked whether the cubic action is the gauge-fixed form of a theory with an independent $2$-string parameter $\Lambda_2$ (\cref{p:sec:gauge}).
  In the classes we have examined, $\Lambda_2$ acts on $\Sigma_2$ alone through a field-independent operator, while $\Psi$ keeps its gauge transformation and $\Sigma_2$ its composite law, on the present space of $2$-string states or on an enlargement of it, including the relaxation of level matching with operators built from the zero modes.
  Every such theory is a trivial enlargement: $\Lambda_2$ shifts $\Sigma_2$ along null directions of the enlarged quadratic form and carries no physics.
  At order $\gs^2$, on the present field content and at generic momentum, the cubic action exposes no other $2$-string symmetry, up to the components that the action does not constrain at this order.
  We do not exclude a theory of which the cubic action is a gauge-fixed form and in which $\Lambda_2$ acts on $\Psi$, or through a field-dependent operator, or through the quartic vertex off level matching.
  On the present field content, a $2$-string symmetry could still appear at order $\gs^3$, which we have not examined.

